Vereinfachend wirke das Drehmoment des Motors direkt (ohne Getriebe) auf die Antriebsräder; Roll- und Luftwiderstand werden vernachlässigt. Die Räder stossen sich am Boden ab: Die Strasse übt am Umfang der Räder, im Abstand $r$ von der Achse, die Antriebskraft
\begin{align}
 F &= \FAuF = \frac{\MAu}{\rAu} = \FAu
\end{align}
auf das Auto aus. Nach dem Aktionsprinzip ist die Beschleunigung
\begin{align}
 a &= \frac{F}{m} = \aAuF = \frac{\MAu}{\rAu\cdot\mAu} = \aAu \approx \result{\aAuP{0}{2}}
\end{align}
